\documentclass[11pt]{article}
\usepackage[T1]{fontenc}
\usepackage{textcomp}
\usepackage[margin=0.75in]{geometry}
\usepackage{amsmath,amssymb,amsthm}
\usepackage{array,booktabs}
\usepackage{hyperref}
\setlength{\parindent}{0pt}
\newtheorem{theorem}{Theorem}
\theoremstyle{definition}
\newtheorem{definition}{Definition}
\title{Flow Proof Artifact: triangle-angle-sum}
\author{Generated by \texttt{flow doc proof}}
\date{Flow Proof Book}
\begin{document}
\maketitle
The interior angles of a triangle sum to two right angles.\\[0.5em]
\textbf{Source.} euclid — Elements, Book I, Proposition 32\\[0.5em]
\bigskip
\noindent{\large \textbf{Derived fact 1.} \textit{The interior angles of a triangle sum to two right angles} \hfill the Euclidean plane \cdot triangle \cdot interior angles sum to two right angles}\par
\smallskip\noindent\textit{Coordinate: the Euclidean plane \cdot triangle \cdot interior angles sum to two right angles}\par
\smallskip\noindent\textit{Source: euclid — Elements, Book I, Proposition 32}\par
\smallskip\noindent\textit{Built on: alternate interior angles are equal when parallel lines meet a transversal}\par
\medskip
\noindent\textbf{Goal.} The interior angles of a triangle sum to two right angles\par
\noindent$\displaystyle \angle A + \angle B + \angle C = 180^\circ$\par
\medskip
\noindent
\begin{tabular}{@{} >{\raggedright\arraybackslash}p{0.06\textwidth} >{\raggedright\arraybackslash}p{0.47\textwidth} >{\raggedleft\arraybackslash}p{0.41\textwidth} @{}}
\textbf{\#} & \textbf{Proof} & \textbf{Mathematics} \\
\midrule
\textbf{1.} & We prove that interior angles sum to two right angles for triangles in the Euclidean plane. & \\[0.45em]
\textbf{2.} & We invoke the axiom governing parallel lines in the Euclidean plane: alternate interior angles are equal when parallel lines meet a transversal. & \\[0.45em]
\textbf{3.} & From step 2, a line through the apex parallel to the base makes alternate angles with the sides, so the three interior angles line up on a straight line, so angle A plus angle B plus angle C equals two right angles. Hence proven. & $\displaystyle \angle A + \angle B + \angle C = 180^\circ$ \\[0.45em]
\bottomrule
\end{tabular}
\label{thm:Geometry-triangle-interior_angles_sum_to_two_right_angles}
\par\medskip\hrule
\end{document}
